\chapter{Preface}
\label{ch:preface}

This book is the first teaching surface for a subject that did not have a name
until the laboratory forced one.

Universities already teach \emph{reconfigurable computing}: LUTs, interconnect,
place-and-route, FPGA fabrics, high-level synthesis.
They already teach \emph{applied cryptography}: AES, RSA, lattices, sometimes
a lecture on fully homomorphic encryption.
They almost never teach what happens when the reconfigurable fabric \emph{is}
a machine-learning graph API, the bitstream is a Yosys netlist, and the values
on the wires are torus samples that must remain encrypted while a host clock
ticks flip-flops.

That intersection is the subject of this course.
We call the working stack \textbf{HELUT}---Homomorphic Edge Look-Up Tensors---and
we call the subject \textbf{reconfigurable homomorphic computing}.

The stack has three pillars~\cite{helut-release,helut-paper}:

\begin{enumerate}
  \item \textbf{Netlist-clocked torus FHE.}
        Gate-level sequential circuits evaluate as exact
        $\mathbb{Z}/2^{32}\mathbb{Z}$ tensor graphs, with real bootstrapping-key
        blind rotation and machine-checkable certificates.
  \item \textbf{Differentiable hardware cryptanalysis.}
        Discrete INIT tables melt into continuous tensors; a $\lambda$-squeeze
        grades shatter versus hold; reciprocal structure that combinatoric
        sieves miss can reappear as an involution.
  \item \textbf{Adversarial polymorphic ciphers.}
        Blue ciphers co-evolve under Red pressure and fail closed, rather than
        waiting to become the next historically leaked stream cipher.
\end{enumerate}

\paragraph{Why a living textbook.}
Research papers freeze a slice of a laboratory.
This course cannot.
The claim inventory~\cite{helut-claim-sheet} moves: Metal kernels get faster,
hedges close, avenues stay unlabeled until they earn receipts.
A professor who taught from a PDF dated June would be teaching a different
Metal compiler than a professor teaching from August.
So the book is versioned against the corpus \emph{epoch}
(\livingepoch{} in this edition) and is honest about stubs.

\paragraph{What this edition is.}
Edition~\corpusedition{} is the \emph{beginnings}: enough foundations to open a
course, enough Pillar~I to run laboratories, enough Pillars~II and~III to
assign readings, and a frontier part that is labeled as frontier.
It is not a complete treatise.
Where the laboratory is thin, the page says so.

\paragraph{What this edition is not.}
It is not a claim that HELUT invented TFHE~\cite{chillotti2020tfhe,ducas2015fhew}.
It is not a production key-size manual until a lattice estimator fills the
pending table (\hid{1}).
It is not a decrypt of U-534 / P1030680.
It is not a side-channel paper.
Those sentences are printed again in Chapter~\ref{ch:living} because students
will forget them the week before the final.

\paragraph{Who wrote it.}
The HELUT project at Digital Defiance.
The scientific voice of the book is the same voice as the research-release
doctrine: lead with what runs, and with what it does not prove.

\paragraph{How to cite an epoch.}
Cite the git commit and the claim-sheet epoch, not ``the HELUT textbook'' as if
it were a static ISBN.
When a student reproduces \cid{18}, they are reproducing a dated Metal NTT
persist tile, not an eternal constant of nature.
